\\documentclass{article}
\\usepackage[margin=1in]{geometry}
\\usepackage{amsmath}
\\usepackage{amssymb}
\\usepackage{mathtools}
\\usepackage{hyperref}
\\hypersetup{hidelinks}
\\numberwithin{equation}{section}
\\begin{document}

\\title{Derivaciones: Memoria por Histéresis y Integración Simpĺectica Implícita}
\\author{Joaqu\\'in St\\'urtz}
\\date{Febrero 2026}
\\maketitle

\\section{Ecuaciones del Flujo Geod\\'esico con Memoria}
Consideramos un estado en espacio de fases $(x, v) \\in \\mathcal{M} \\times T\\mathcal{M}$ que evoluciona bajo una ecuaci\\'on geod\\'esica forzada con disipaci\\'on y un t\\'ermino de fuerza ``fantasma'' (memoria por hist\\'eresis):
\\begin{align}
\\dot{x} &= v,\\\\
\\dot{v} &= F_{\\text{ext}}(u) - \\Gamma(x, v) - \\mu(x, u)\\,v + F_{\\text{ghost}}(h),
\\end{align}
donde $\\Gamma$ es la resistencia geom\\'etrica (Christoffel efectiva), $\\mu \\ge 0$ es el coeficiente de fricci\\'on selectiva, y $h$ es el estado de memoria por hist\\'eresis.

\\subsection{Memoria de Hist\\'eresis con Complejidad $\\mathcal{O}(1)$}
Definimos el estado de memoria $h \\in \\mathbb{R}^d$ (misma dimensionalidad que $x, v$) y su actualizaci\\'on discreta:
\\begin{align}
h_{t+1} &= \\lambda\\,h_t + \\tanh\\!\\big( W_u\\,\\phi(x_t) + W_v\\,v_t + b \\big), \\quad 0 < \\lambda < 1,\\\\
F_{\\text{ghost}}(h_t) &= W_g\\,h_t,
\\end{align}
donde $\\phi(x) = [\\sin(x), \\cos(x)]$ en topolog\\'ias toroidales y $\\phi(x) = x$ en variedades euclidianas. $W_u, W_v, W_g$ son matrices aprendibles y $b$ un sesgo. El estado $h_t$ es de tama\\~no fijo independiente de la longitud de secuencia $L$, por lo que la memoria de inferencia es $\\mathcal{O}(1)$.

\\paragraph{Correcci\\'on recomendada del t\\'ermino de memoria.}
Para evitar saturaci\\'on prematura y pozos sem\\'anticos excesivos, sustituir $\\tanh$ por $\\mathrm{Softsign}$ o $\\tanh$ con ganancia moderada:
\\begin{align}
\\mathrm{Softsign}(z) &= \\frac{z}{1 + |z|}, \\quad \\text{o}\\quad \\tanh(\\alpha z),\\ \\alpha \\in [0.5, 1.5].
\\end{align}
Esto mantiene la estabilidad de $h$ y reduce la sensibilidad a picos locales.

\\section{Paso Simp\\'lectico con Fricci\\'on Impl\\'icita}
Usamos Leapfrog (kick--drift--kick) con fricci\\'on impl\\'icita y memoria fantasma. Sea $\\Delta t$ el paso base y $h = \\tfrac{1}{2}\\,g(x)\\,\\Delta t$ con compresi\\'on temporal $g(x) \\in (0,1]$:
\\begin{align}
v_{t+\\frac{1}{2}} &= \\frac{v_t + h\\,\\big(F_{\\text{ext}}(u_t) + F_{\\text{ghost}}(h_t) - \\Gamma(x_t, v_t)\\big)}{1 + h\\,\\mu(x_t, u_t)},\\\\
x_{t+1} &= x_t + \\Delta t_{\\text{eff}}\\,v_{t+\\frac{1}{2}},\\quad \\Delta t_{\\text{eff}} = g(x_t)\\,\\Delta t,\\\\
v_{t+1} &= \\frac{v_{t+\\frac{1}{2}} + h\\,\\big(F_{\\text{ext}}(u_t) + F_{\\text{ghost}}(h_t) - \\Gamma(x_{t+1}, v_{t+\\frac{1}{2}})\\big)}{1 + h\\,\\mu(x_{t+1}, u_t)}.
\\end{align}
Este esquema conserva volumen de fase para $\\mu=0$ y controla la irreversibilidad mediante $\\mu>0$.

\\subsection{Cota de Energ\\'ia y Estabilidad en Float32}
La energ\\'ia cin\\'etica $E_k = \\tfrac{1}{2}\\,\\|v\\|^2$ satisface, bajo fricci\\'on positiva,
\\begin{align}
E_k(t+1) \\le \\frac{E_k(t)}{\\big(1 + h\\,\\mu_{\\min}\\big)^2} + C\\,\\Delta t\\,\\|F_{\\text{ext}} + F_{\\text{ghost}}\\|,
\\end{align}
para alguna constante $C$ dependiente de curvatura y acoplamientos locales. La divisi\\'on impl\\'icita por $1 + h\\,\\mu$ act\\'ua como contracci\\'on y aten\\'ua drift num\\'erico en $\\mathrm{float32}$ siempre que $\\mu_{\\min} > 0$ en regiones de alta curvatura.

\\section{Crist\\'offel Reactivo con Clamping}
El Christoffel efectivo incluye modulaci\\'on reactiva y protecci\\'on de divisiones:
\\begin{align}
\\Gamma_{\\text{eff}}(x, v) &= \\Gamma_{\\text{base}}(x, v)\\Big[1 + \\rho\\big(\\mathcal{E}(v)\\big)\\Big],\\\\
\\rho(\\mathcal{E}) &= \\mathrm{clip}\\big(\\alpha\\,\\tanh(\\mathcal{E}),\\ 0,\\ \\rho_{\\max}\\big),\\quad \\mathcal{E}(v) = \\tanh\\big(\\|v\\|^2\\big),
\\end{align}
y en toros:
\\begin{align}
\\Gamma^{\\phi}_{\\text{torus}}(x) &= -\\,\\frac{r\\,\\sin(\\theta)}{\\max\\big(R + r\\,\\cos\\theta,\\ \\epsilon_{\\text{clamp}}\\big)},\\quad \\epsilon_{\\text{clamp}} \\approx 10^{-1}.
\\end{align}
La doble protecci\\'on (clamp + $\\epsilon$) evita divisiones cercanas a cero y aliasing en pasos largos.

\\section{Memoria $\\mathcal{O}(1)$ y Complejidad}
El sistema mantiene un estado compacto $(x, v, h)$ de tama\\~no fijo, por lo que:
\\begin{align}
\\text{Memoria de inferencia} &= \\mathcal{O}(1),\\\\
\\text{Tiempo por secuencia} &\\approx \\mathcal{O}(L\\,d^2)\\ \\text{(mezcla densa por capa)}.
\\end{align}
La memoria no crece con $L$; $h$ es un acumulador geom\\'etrico interno que reemplaza caches externos.

\\section{Observaciones Pr\\'acticas}
\\begin{itemize}
\\item Ajustar $\\lambda \\in [0.85, 0.95]$ y $\\alpha \\in [0.5, 1.0]$ para evitar atractores sem\\'anticos permanentes.
\\item Usar $g(x)$ para contraer tiempo en regiones de alta curvatura y prevenir aliasing.
\\item Verificar drift de energ\\'ia $<1\\%$ en 10k pasos con gates activos.
\\end{itemize}

\\end{document}
